\centerline{\bf \Huge Оптимизация}
\centerline{\bf \Huge }


\problem{\textbf{1}} В прямоугольной комнате площадью 42 м$^2$ требуется установить плинтусы по всему периметру. Стоимость 1 метра плинтуса составляет 280 рублей. При каких целых линейных размерах
комнаты затраты на покупку плинтуса будут наименьшими?

\problem{\textbf{2}} Компания изготавливает и продает изделия. Если одно изделие стоит 2000 рублей, то реализуется 1000 штук изделий. При снижении средней цены одного изделия на 50 рублей
объемы реализации возрастают на 50 штук. При какой цене фирма получит максимальный
доход и каково его значение?

\problem{\textbf{3}} На двух заводах, которыми владеет Александр, производят одинаковый товар. Если на первом заводе рабочие суммарно трудятся $t^2$
 часов в неделю, то они производят $t$ товаров. Если
на втором заводе рабочие трудятся $t^2$
часов в неделю, то они производят $2t$ товаров. Заработная
плата рабочего за час работы составляет 300 рублей. Найдите наименьшую сумму, которую
должен потратить на зарплаты рабочим в неделю Александр, чтобы оба завода произвели 600
единиц товара. Ответ дайте в млн. рублей.

\problem{\textbf{4}} Предприниматель купил здание и собирается открыть в нём отель. В отеле могут быть
стандартные номера площадью 30 квадратных метров и номера «люкс» площадью 40 квадратных метров. Общая площадь, которую можно отвести под номера, составляет 940 квадратных
метров. Предприниматель может поделить эту площадь между номерами различных типов,
как хочет. Обычный номер будет приносить отелю 4000 рублей в сутки, а номер «люкс» — 5000
рублей в сутки. Какую наибольшую сумму денег сможет заработать в сутки на своём отеле
предприниматель?


\problem{\textbf{5}} Строительство нового завода стоит 78 млн рублей. Затраты на производство $x$ тыс. единиц продукции на таком заводе равны $0,5x^2+ 2x + 6$ млн рублей в год. Если продукцию завода
продать по цене $p$ тыс. рублей за единицу, то прибыль фирмы (в млн рублей) за один год
составит $px - (0,5x^2+ 2x + 6)$.
Когда завод будет построен, фирма будет выпускать продукцию в таком количестве, чтобы
прибыль была наибольшей.
При каком наименьшем значении $p$ строительство завода окупится не более чем за 3 года?

\problem{\textbf{6}} В двух областях есть по 20 рабочих, каждый из которых готов трудиться по 10 часов в сутки на добыче алюминия или никеля. В первой области один рабочий за час добывает 0,2 кг
алюминия или 0,2 кг никеля. Во второй области для добычи $x$ кг алюминия в день требуется
$x^2$ человеко-часов труда, а для добычи $y$ кг никеля в день требуется $y^2$ человеко-часов труда.
Обе области поставляют добытый металл на завод, где для нужд промышленности производится сплав алюминия и никеля, в котором на 1 кг алюминия приходится 1 кг никеля. При этом
области договариваются между собой вести добычу металлов так, чтобы завод мог произвести
наибольшее количество сплава. Сколько килограммов сплава при таких условиях ежедневно
сможет произвести завод?

\problem{\textbf{7}} Фирма имеет возможность рекламировать свою продукцию, используя местных радио и телевизионную сети. Затраты на рекламу в бюджете фирмы ограничены величиной 1000\$ в
месяц. Каждая минута радиорекламы обходится в 5\$, а каждая минута телерекламы – в 100\$.
Фирма хотела бы использовать радиосеть, по крайней мере, в два раза чаще, чем сеть телевидения, но при этом фирма решила, что время радиорекламы не должно превышать двух
часов. Опыт прошлых лет показал, что объем сбыта, который обеспечивает каждая минута телерекламы, в 25 раз больше сбыта, обеспечиваемого одной минутой радиорекламы. Определите
оптимальное распределение финансовых средств, ежемесячно отпускаемых на рекламу, между
радио и телерекламой, если время можно покупать только поминутно.

\problem{\textbf{8}} Зависимость количества $Q$ в шт. при условии $0 \leqslant Q \leqslant 15000$ купленного у фирмы товара
от цены $P$ в руб. за шт. выражается формулой $Q$ = 15000 - $P$. Затраты на производство $Q$ единиц
товара составляют $3000Q$ + 1000000 рублей. Кроме затрат на производство, фирма должна
платить налог $t$ рублей при условии 0 < $t$ < 10000 с каждой произведённой единицы товара.
Таким образом, прибыль фирмы составляет $PQ - 3000Q - 1000000 - tQ$ рублей, а общая сумма
налогов, собранных государством, равна $tQ$ рублей.
Фирма производит такое количество товара, при котором её прибыль максимальна. При
каком значении $t$ общая сумма налогов, собранных государством, будет максимальной?

\problem{\textbf{9}} В двух шахтах добывают алюминий и никель. В первой шахте имеется 60 рабочих, каждый из которых готов трудиться 5 часов в день. При этом один рабочий за час добывает 2 кг
алюминия или 3 кг никеля. Во второй шахте имеется 260 рабочих, каждый из которых готов
трудиться 5 часов в день. При этом один рабочий за час добывает 3 кг алюминия или 2 кг
никеля.
Обе шахты поставляют добытый металл на завод, где для нужд промышленности производится сплав алюминия и никеля, в котором на 2 кг алюминия приходится 1 кг никеля. $MAX$ кг сплава?

